\documentclass[10pt,a4paper]{amsart}
\usepackage{amsmath,amssymb,amsthm,graphicx}
% Retained minimal preamble for the Stacks Project chapter varieties.tex.
% Verbatim excerpt of the chapter preamble: Xy-pic setup (native preamble lines 14-18), the native theorem
% declarations with the mathematical macros (native preamble lines 154-217), and the project own no-op
% definitions of the reference/slogan/history display environments (native wrapper.tex lines 1-3) that the
% counted statement uses for its slogan. Copied mechanically; nothing added, removed or redefined.
% For commutative diagrams we use Xy-pic
\usepackage[all]{xy}

% We use 2cell for 2-commutative diagrams.
\xyoption{2cell}
% Theorem environments.
%
\theoremstyle{plain}
\newtheorem{theorem}[subsection]{Theorem}
\newtheorem{proposition}[subsection]{Proposition}
\newtheorem{lemma}[subsection]{Lemma}

\theoremstyle{definition}
\newtheorem{definition}[subsection]{Definition}
\newtheorem{example}[subsection]{Example}
\newtheorem{exercise}[subsection]{Exercise}
\newtheorem{situation}[subsection]{Situation}

\theoremstyle{remark}
\newtheorem{remark}[subsection]{Remark}
\newtheorem{remarks}[subsection]{Remarks}

\numberwithin{equation}{subsection}

% Macros
%
\def\lim{\mathop{\mathrm{lim}}\nolimits}
\def\colim{\mathop{\mathrm{colim}}\nolimits}
\def\Spec{\mathop{\mathrm{Spec}}}
\def\Hom{\mathop{\mathrm{Hom}}\nolimits}
\def\Ext{\mathop{\mathrm{Ext}}\nolimits}
\def\SheafHom{\mathop{\mathcal{H}\!\mathit{om}}\nolimits}
\def\SheafExt{\mathop{\mathcal{E}\!\mathit{xt}}\nolimits}
\def\Sch{\mathit{Sch}}
\def\Mor{\mathop{\mathrm{Mor}}\nolimits}
\def\Ob{\mathop{\mathrm{Ob}}\nolimits}
\def\Sh{\mathop{\mathit{Sh}}\nolimits}
\def\NL{\mathop{N\!L}\nolimits}
\def\CH{\mathop{\mathrm{CH}}\nolimits}
\def\proetale{{pro\text{-}\acute{e}tale}}
\def\etale{{\acute{e}tale}}
\def\QCoh{\mathit{QCoh}}
\def\Ker{\mathop{\mathrm{Ker}}}
\def\Im{\mathop{\mathrm{Im}}}
\def\Coker{\mathop{\mathrm{Coker}}}
\def\Coim{\mathop{\mathrm{Coim}}}

% Boxtimes
%
\DeclareMathSymbol{\boxtimes}{\mathbin}{AMSa}{"02}

%
% Macros for moduli stacks/spaces
%
\def\QCohstack{\mathcal{QC}\!\mathit{oh}}
\def\Cohstack{\mathcal{C}\!\mathit{oh}}
\def\Spacesstack{\mathcal{S}\!\mathit{paces}}
\def\Quotfunctor{\mathrm{Quot}}
\def\Hilbfunctor{\mathrm{Hilb}}
\def\Curvesstack{\mathcal{C}\!\mathit{urves}}
\def\Polarizedstack{\mathcal{P}\!\mathit{olarized}}
\def\Complexesstack{\mathcal{C}\!\mathit{omplexes}}
% \Pic is the operator that assigns to X its picard group, usage \Pic(X)
% \Picardstack_{X/B} denotes the Picard stack of X over B
% \Picardfunctor_{X/B} denotes the Picard functor of X over B
\def\Pic{\mathop{\mathrm{Pic}}\nolimits}
\def\Picardstack{\mathcal{P}\!\mathit{ic}}
\def\Picardfunctor{\mathrm{Pic}}
\def\Deformationcategory{\mathcal{D}\!\mathit{ef}}
\newenvironment{reference}{}{}
\newenvironment{slogan}{}{}
\newenvironment{history}{}{}
\usepackage{cleveref}
\crefname{originaldefinition}{original definition}{original definitions}
\Crefname{originaldefinition}{Original definition}{Original definitions}
\makeatletter
\newlabel{algebra-definition-geometrically-regular}{{0382}{source}{}{sourcepointer-1}{}}
\newlabel{algebra-definition-geometrically-regular@cref}{{[originaldefinition][0][]0382}{[1][0][]source}{}{}{}}
\makeatother
\crefname{originallemma}{original lemma}{original lemmas}
\Crefname{originallemma}{Original lemma}{Original lemmas}
\makeatletter
\newlabel{properties-lemma-characterize-regular}{{02IT}{source}{}{sourcepointer-2}{}}
\newlabel{properties-lemma-characterize-regular@cref}{{[originallemma][0][]02IT}{[1][0][]source}{}{}{}}
\makeatother
\crefname{originalsection}{original section}{original sections}
\Crefname{originalsection}{Original section}{Original sections}
\makeatletter
\newlabel{schemes-section-fibre-products}{{01JO}{source}{}{sourcepointer-3}{}}
\newlabel{schemes-section-fibre-products@cref}{{[originalsection][0][]01JO}{[1][0][]source}{}{}{}}
\makeatother
\title{Five equivalent characterizations of geometric regularity for a locally Noetherian scheme over a field}
\author{The Stacks Project Authors}
\date{stacks-project:varieties.tex:local-snapshot-20261011; GNU FDL 1.2 or later}
\begin{document}
\maketitle
\noindent\textit{Editorial note.} All delivered mathematical text is the author's own, copied from the pinned Stacks Project chapter file varieties.tex; no statement or proof step was rewritten, reordered or removed.\par\medskip
\noindent\textit{Editorial note.} The proof is complete as written and closes the cycle of implications; no step is referred to the reader.\par\medskip
\noindent\textit{Editorial note.} The counted spans carry 3 navigation stubs and no bibliography citation.\par\medskip
\noindent\textit{Editorial note.} The preamble is the chapter's retained minimal preamble (sources/stacks-varieties/retained-preamble.tex lines 1-77), a verbatim excerpt of the native preamble.\par\medskip
\noindent\textit{Editorial note.} The label lemma-geometrically-regular is defined only in varieties.tex, so no cross-chapter restatement had to be weighed for this card.\par\medskip
\noindent\textit{Allowed standard prerequisites (editorial).}\begin{itemize}
\item Criterion for a scheme to be regular through its local rings. Properties, Lemma properties-lemma-characterize-regular, invoked at proof lines 2446-2447. Supplies the transition from regular local rings to a regular scheme in the first implication.
\item Geometric regularity of a k-algebra. Algebra, Definition algebra-definition-geometrically-regular, named in the statement and used at proof lines 2457-2459. Supplies the algebraic property to which conditions (3) and (4) are translated.
\item Base change as gluing of spectra of tensor products of coordinate rings. Schemes, Section schemes-section-fibre-products, invoked at proof lines 2463-2466. Supplies the final implication from the coordinate rings of one affine covering to geometric regularity.
\end{itemize}\medskip
\ifdefined\hypertarget\hypertarget{sourcepointer-1}{}\fi
\noindent\textit{Source reference.} Original definition 0382: sources/stacks-chow/tags.txt, source lines 4221--4221. Pointer outside the counted span, invoked once for context; the load-bearing argument of the counted statement is the author's proof reproduced in full in the proof span.\par\medskip
\ifdefined\hypertarget\hypertarget{sourcepointer-2}{}\fi
\noindent\textit{Source reference.} Original lemma 02IT: sources/stacks-chow/tags.txt, source lines 3312--3312. Pointer outside the counted span, invoked once for context; the load-bearing argument of the counted statement is the author's proof reproduced in full in the proof span.\par\medskip
\ifdefined\hypertarget\hypertarget{sourcepointer-3}{}\fi
\noindent\textit{Source reference.} Original section 01JO: sources/stacks-chow/tags.txt, source lines 2021--2021. Pointer outside the counted span, invoked once for context; the load-bearing argument of the counted statement is the author's proof reproduced in full in the proof span.\par\medskip
\begin{lemma}
\label{lemma-geometrically-regular}
Let $k$ be a field.
Let $X$ be a locally Noetherian scheme over $k$.
The following are equivalent
\begin{enumerate}
\item $X$ is geometrically regular,
\item $X_{k'}$ is a regular scheme for every finitely generated field
extension $k'/k$,
\item $X_{k'}$ is a regular scheme for every finite purely inseparable
field extension $k'/k$,
\item for every affine open $U \subset X$ the ring $\mathcal{O}_X(U)$
is geometrically regular (see
Algebra, Definition \ref{algebra-definition-geometrically-regular}), and
\item there exists an affine open covering $X = \bigcup U_i$ such that
each $\mathcal{O}_X(U_i)$ is geometrically regular over $k$.
\end{enumerate}
\end{lemma}

\begin{proof}
Assume (1). Then for every finitely generated field extension
$k'/k$ and
every point $x' \in X_{k'}$ the local ring of $X_{k'}$ at $x'$
is regular. By Properties, Lemma \ref{properties-lemma-characterize-regular}
this means that $X_{k'}$ is regular. Hence (2).

\medskip\noindent
It is clear that (2) implies (3).

\medskip\noindent
Assume (3) and let $U \subset X$ be an affine open subscheme.
Then $U_{k'}$ is a regular scheme for any finite purely inseparable
extension $k'/k$ (including $k = k'$). This means that
$k' \otimes_k \mathcal{O}(U)$ is a regular ring for all
finite purely inseparable extensions $k'/k$. Hence
$\mathcal{O}(U)$ is a geometrically regular $k$-algebra
and we see that (4) holds.

\medskip\noindent
It is clear that (4) implies (5). Let $X = \bigcup U_i$ be an affine
open covering as in (5). For any field extension $k'/k$ the base
change $X_{k'}$ is gotten by gluing the spectra of the
rings $\mathcal{O}_X(U_i) \otimes_k k'$ (see
Schemes, Section \ref{schemes-section-fibre-products}).
Hence $X_{k'}$ is regular. So (1) holds.
\end{proof}

\end{document}
